% =====================================================================
%  Thème 5 — EXERCICES — Niveau MOYEN
% =====================================================================
\documentclass[11pt,a4paper]{article}
\usepackage[utf8]{inputenc}
\usepackage[T1]{fontenc}
\usepackage{lmodern}
\usepackage[french]{babel}
\usepackage{amsmath,amssymb,mathtools}
\usepackage{icomma}
\usepackage{xcolor}
\usepackage[most]{tcolorbox}
\usepackage{enumitem}
\usepackage{array}
\usepackage{booktabs}
\usepackage{fancyhdr}
\usepackage[hidelinks]{hyperref}
\usepackage[a4paper,margin=2.2cm,headheight=26pt,headsep=10pt]{geometry}

\definecolor{primary}{RGB}{223,87,99}
\definecolor{primarydark}{RGB}{150,42,55}

\tcbset{boxbase/.style={enhanced,breakable,boxrule=0.9pt,arc=2.5pt,
   left=3mm,right=3mm,top=2.3mm,bottom=2.3mm,
   fonttitle=\bfseries,coltitle=white,toptitle=0.6mm,bottomtitle=0.6mm}}
\newtcolorbox[auto counter]{exercice}[1][]{boxbase,colback=primary!4,colframe=primary,
  title={\textsc{Exercice}~\thetcbcounter\ifx\relax#1\relax\relax\else\textmd{ --- \itshape #1}\fi}}

\pagestyle{fancy}
\fancyhf{}
\fancyhead[L]{\small\textcolor{primarydark}{\textbf{Fiche 6} — Les limites}}
\fancyhead[R]{\small\textcolor{primarydark}{\textsc{Thème 5 — Moyen}}}
\fancyfoot[C]{\small\thepage}
\renewcommand{\headrulewidth}{0.8pt}
\renewcommand{\headrule}{\hbox to\headwidth{\color{primary}\leaders\hrule height \headrulewidth\hfill}}

\newcommand{\R}{\mathbb{R}}

\begin{document}

\begin{center}
{\Large\bfseries\color{primarydark}Thème 5 — L'Hospital \& limites trigonométriques}\\[2pt]
{\large\color{primary}Exercices — Niveau Moyen}
\end{center}
\vspace{4pt}
{\small\itshape On combine l'Hospital, les limites remarquables, et la transformation $0\times\infty\to$
quotient.}
\vspace{6pt}

\begin{exercice}[l'Hospital avec le cosinus]
Calculer $\displaystyle\lim_{x\to \frac{\pi}{2}}\frac{\cos x}{x-\frac{\pi}{2}}$.
\begin{enumerate}[label=\alph*),leftmargin=2em,itemsep=3pt]
  \item Vérifier que la forme est $\frac00$.
  \item Appliquer l'Hospital et conclure.
\end{enumerate}
\end{exercice}

\begin{exercice}[combinaison de limites remarquables]
Calculer $\displaystyle\lim_{x\to 0}\frac{\sin 2x}{\tan x}$ en décomposant en facteurs de référence
$\dfrac{\sin 2x}{2x}$ et $\dfrac{x}{\tan x}$.
\end{exercice}

\begin{exercice}[transformer $0\times\infty$]
Calculer $\displaystyle\lim_{x\to 0^{+}} x\ln x$.
\begin{enumerate}[label=\alph*),leftmargin=2em,itemsep=3pt]
  \item Quelle est la forme directe ? Pourquoi est-elle indéterminée ?
  \item Réécrire $x\ln x=\dfrac{\ln x}{1/x}$, puis appliquer l'Hospital.
\end{enumerate}
\end{exercice}

\begin{exercice}[l'Hospital deux fois]
Calculer $\displaystyle\lim_{x\to 0}\frac{1-\cos x}{x^2}$.
\begin{enumerate}[label=\alph*),leftmargin=2em,itemsep=3pt]
  \item Vérifier la forme $\frac00$, appliquer l'Hospital une première fois.
  \item La forme est encore $\frac00$ : appliquer l'Hospital une seconde fois et conclure.
\end{enumerate}
\end{exercice}

\begin{exercice}[la même expression, deux points]
Soit $q(x)=\dfrac{e^x-1}{x}$.
\begin{enumerate}[label=\alph*),leftmargin=2em,itemsep=3pt]
  \item Calculer $\displaystyle\lim_{x\to 0} q(x)$ (forme $\frac00$).
  \item Calculer $\displaystyle\lim_{x\to+\infty} q(x)$ (forme $\frac{\infty}{\infty}$).
  \item Que traduisent ces deux résultats sur la « vitesse » de croissance de $e^x$ face à $x$ ?
\end{enumerate}
\end{exercice}

\end{document}
